\chapter{Introduction to Weighted Directional Total Nuclear Variation}
\label{chapter:wdtnv}
In this chapter, we introduce a smoothed \acrfull{wdtnv} prior for synergistic emission tomography. The chapter focuses on the regulariser itself: its relationship to total nuclear variation, its differentiability after singular-value smoothing, the use of anatomical information through directional operators, and the modality-wise weighting needed when coupling heterogeneous emission images.

The optimisation machinery used to apply this prior in large-scale \gls{pet}/\gls{spect} reconstruction, including prior-aware preconditioning, subset selection, and multi-bed \gls{pet} handling, is developed separately in Chapter~\ref{chapter:optimisation}.

\section{Total Nuclear Variation Revisited}
\label{sec:tv_to_tnv}
This section builds on the definitions in Section~\ref{subsec:synergistic_recon}. For convenience, we restate the key definitions here:
\begin{itemize}
    \item Registered images: $\bm{z} \coloneqq (z_1,\dots,z_M)\in\mathbb{R}_+^{N\times M}$, where the matrix-space identification stacks the images as columns, $z_m \coloneqq W_m x_m$, and
    $W_m : \mathbb{R}_+^{J_m} \to \mathbb{R}_+^{N}$.
    \item Finite-difference operator: for $\ell=1,\dots,d$,
    \[
        \big((D x)_n\big)^{(\ell)} \coloneqq \frac{x_{n+\delta_\ell} - x_{n}}{\Delta_{\ell}} .
    \]
    \item Voxel-wise Jacobian:
    \[
        \mathcal{G}_n(\bm{z})
        \coloneqq
        \begin{bmatrix}
        (D z_1)_n & (D z_2)_n & \cdots & (D z_M)_n
        \end{bmatrix}
        \in \mathbb{R}^{d\times M}.
    \]
    \item Nuclear norm: for a matrix $Y\in\mathbb{R}^{d\times M}$ with singular values $\{\sigma_k(Y)\}_{k=1}^{r}$ and $r=\min\{d,M\}$,
    \[
        \|Y\|_* \coloneqq \sum_{k=1}^{r}\sigma_k(Y).
    \]
    \item Total nuclear variation:
    \[
        \mathcal{R}_{\mathrm{TNV}}(\bm{z}) \coloneqq \sum_{n=1}^N  \|\mathcal{G}_n(\bm{z})\|_* .
    \]
\end{itemize}

\subsection{Mathematical properties and differentiability}
\label{subsec:tnv_differentiability}

Briefly revisiting total variation, $\mathcal{R}_{\mathrm{TV}}$ (cf.\ \eqref{eq:tv}) is not differentiable when a voxel gradient $(D x_m)_j$ vanishes. For gradient-based optimisation, we use a Charbonnier smoothing with $\epsilon>0$.

If we define
\begin{equation}
\tilde{\mathcal{R}}_{\mathrm{TV}}(x_m)
\;:=\;
\sum_{j=1}^{J_m} \|(D x_m)_j\|_{2,\epsilon},
\qquad
\|v\|_{2,\epsilon}:=\sqrt{\|v\|_2^2+\epsilon^2}-\epsilon,
\end{equation}
then the image-space gradient is
\begin{equation}
\nabla_{x_m}\tilde{\mathcal{R}}_{\mathrm{TV}}(x_m)
\;=\;
-\,\operatorname{div}\!\left(\frac{D x_m}{\sqrt{\|D x_m\|_2^2+\epsilon^2}}\right),
\end{equation}
where the division is understood pointwise over voxels.

The nuclear norm used by \gls{tnv} is a particular orthogonally invariant penalty: $\|UYV^\top\|_*=\|Y\|_*$ for orthogonal $U\in\mathbb{R}^{d\times d}$ and $V\in\mathbb{R}^{M\times M}$. More generally, an orthogonally invariant scalar function of $Y\in\mathbb{R}^{d\times M}$ can be written as $F(Y)=f(\sigma(Y))$, where $\sigma(Y)$ lists its $r=\min\{d,M\}$ singular values and $f$ is \emph{absolutely symmetric}, i.e. it is invariant under permutations and sign changes of its arguments \cite{lewisConvexAnalysisUnitarily1995}.  This is obviously the case for a sum of non-negative singular values.

We can smooth TNV by applying the same Charbonnier function used for \gls{tv} to each singular value \cite{charbonnier_deterministic_1997,fessler_image_2017}:
\begin{equation}\label{eq:charbonnier_g}
g(s)=\sqrt{s^2+\epsilon^2}-\epsilon,
\qquad
g'(s)=\frac{s}{\sqrt{s^2+\epsilon^2}},
\qquad
g''(s)=\frac{\epsilon^2}{(s^2+\epsilon^2)^{3/2}},
\end{equation}
where $\epsilon>0$ controls the degree of smoothing.

For the smoothed penalty, define
\begin{equation}
    \phi(Y):=f_\epsilon(\sigma(Y)),
    \qquad
    f_\epsilon(s):=\sum_{\ell=1}^{r}g(s_\ell).
\end{equation}
The Charbonnier function $g$ in \eqref{eq:charbonnier_g} is even, convex and smooth across zero, so $f_\epsilon$ is convex, absolutely symmetric and differentiable on all of $\mathbb{R}^r$. Lewis's gradient formula for convex orthogonally invariant functions therefore shows that $\phi$ is differentiable even when $Y$ is rank deficient \cite[Theorem~3.1]{lewisConvexAnalysisUnitarily1995}. This is the matrix analogue of smoothing the corner of the absolute value in scalar \gls{tv}.


\paragraph{Convexity.}
Since $f_\epsilon$ is convex and absolutely symmetric, $\phi=f_\epsilon\circ\sigma$ is convex on $\mathbb{R}^{d\times M}$ \cite[Corollary~2.6]{lewisConvexAnalysisUnitarily1995}. The finite-difference operator $D$, directional operators $\Pi_{m,n}$ of Section~\ref{subsec:directional_from_guidance}, and weight matrices $B_n$ of Section~\ref{subsec:modality_weighting} are fixed linear maps. Composing $\phi$ with these maps and summing over voxels preserves convexity, so
\begin{equation}
    \mathcal{R}_{\mathrm{w\text{-}dTNV}} \ \text{is convex on}\ \mathbb{R}^{N\times M}.
    \label{eq:wdtnv_convex}
\end{equation}

\paragraph{Jacobian-space gradient.}
For a voxel-wise Jacobian $\mathcal{G}\in\mathbb{R}^{d\times M}$, a thin \gls{svd}
\begin{equation}
\mathcal{G} \;=\; U\,\Sigma\,V^\top,
\qquad
U\in\mathbb{R}^{d\times r},\ \Sigma=\mathrm{diag}(\sigma)\in\mathbb{R}^{r\times r},\ V\in\mathbb{R}^{M\times r},
\end{equation}
with $r=\min\{d,M\}$ gives
\begin{equation}\label{eq:phi_grad_svd}
\nabla_{\mathcal{G}} \phi(\mathcal{G})
\;=\;
U\,\mathrm{diag}\!\bigl(g'(\sigma)\bigr)\,V^\top,
\qquad
g'(\sigma_\ell)=\frac{\sigma_\ell}{\sqrt{\sigma_\ell^2+\epsilon^2}}.
\end{equation}

In practice, it is unnecessary to form the left singular vectors, $U$, if we exploit the small Gram matrix
\begin{equation}
G := \mathcal{G}^\top \mathcal{G}\in\mathbb{R}^{M\times M},
\qquad G\succeq 0.
\end{equation}
The penalty itself is a function of the eigenvalues of this symmetric matrix:
\begin{equation}
    \phi(\mathcal{G})
    =\sum_{i=1}^{M}\left(\sqrt{\lambda_i(G)+\epsilon^2}-\epsilon\right)
    =\operatorname{tr}\!\left((G+\epsilon^2\mathbb{I}_M)^{1/2}\right)-M\epsilon.
    \label{eq:phi_gram_spectral}
\end{equation}
Any additional zero eigenvalues when $M>d$ contribute zero. This representation also connects the prior to the calculus of spectral functions of symmetric matrices used for its Hessian in Chapter~\ref{chapter:optimisation}.
Now, $G=V\Lambda V^\top$\ is a real-valued Hermitian matrix with $\Lambda=\mathrm{diag}(\lambda_1,\dots,\lambda_M)$. It has eigenvalues $\lambda_\ell=\sigma_\ell^2$ and $V$ coincides with the right singular vectors of $\mathcal{G}$. Substituting $U=\mathcal{G}V\Sigma^{-1}$\footnote{This follows from right-multiplying the singular value decomposition of $\mathcal{G}$ by $V \Sigma^{-1}$.} into~\eqref{eq:phi_grad_svd} gives
\begin{equation}\label{eq:phi_grad_gram}
\nabla_{\mathcal{G}} \phi(\mathcal{G})
\;=\;
\mathcal{G}\,V\,\mathrm{diag}\!\left(\frac{g'(\sigma_\ell)}{\sigma_\ell}\right)\,V^\top,
\end{equation}
with the continuous extension at $\sigma_\ell=0$. For Charbonnier smoothing,
\begin{equation}
\frac{g'(\sigma_\ell)}{\sigma_\ell}
=
\frac{1}{\sqrt{\sigma_\ell^2+\epsilon^2}}
=
\frac{1}{\sqrt{\lambda_\ell+\epsilon^2}},
\end{equation}
hence
\begin{equation}\label{eq:phi_grad_inv_sqrt}
\nabla_{\mathcal{G}} \phi(\mathcal{G})
\;=\;
\mathcal{G}\,V\,\mathrm{diag}\!\left(\frac{1}{\sqrt{\lambda_\ell+\epsilon^2}}\right)\,V^\top
\;=\;
\mathcal{G}\,(G+\epsilon^2 \mathbb{I})^{-1/2}.
\end{equation}
Equation~\eqref{eq:phi_grad_inv_sqrt} shows that the Jacobian-space gradient is obtained by applying the inverse square-root of the $M\times M$ \gls{spd} matrix $G+\epsilon^2 \mathbb{I}$ to the columns of $\mathcal{G}$ \footnote{Note that if $d<M$ the same strategy can be applied to $\mathcal{G}\mathcal{G}^\top$}. This object recurs throughout the remainder of this chapter and in Chapters~\ref{chapter:optimisation} and~\ref{discussion}, so it is worth naming. For any $Y\in\mathbb{R}^{d\times M}$ define the diffusivity matrix
\begin{equation}
    \Omegaeps{Y} \coloneqq \left(Y^\top Y+\epsilon^2\mathbb{I}\right)^{-1/2}
    \in\mathbb{R}^{M\times M},
    \label{eq:spectral_weight}
\end{equation}
so that~\eqref{eq:phi_grad_inv_sqrt} reads $\nabla_{\mathcal{G}}\phi(\mathcal{G})=\mathcal{G}\,\Omegaeps{\mathcal{G}}$. Because $\epsilon>0$, $\Omegaeps{Y}$ is well defined and \gls{spd} for every $Y$, including rank-deficient ones.

        For $M=2$ modalities, we can even avoid an eigendecomposition of $G$. Writing
        \begin{equation}
        G=
        \begin{pmatrix}
        a & c\\
        c & b
        \end{pmatrix},
        \qquad
        a=\|\mathcal{G}_{:,1}\|_2^2,\ \ b=\|\mathcal{G}_{:,2}\|_2^2,\ \ c=\langle \mathcal{G}_{:,1},\mathcal{G}_{:,2}\rangle,
        \end{equation}
        define $S:=G+\epsilon^2 \mathbb{I}=\begin{psmallmatrix}s_{11} & s_{12}\\ s_{12} & s_{22}\end{psmallmatrix}$ with $s_{11}=a+\epsilon^2$, $s_{22}=b+\epsilon^2$, $s_{12}=c$, and the invariants
        \begin{equation}
        t:=\mathrm{tr}(S)=s_{11}+s_{22},
        \qquad
        s:=\sqrt{\det(S)}=\sqrt{s_{11}s_{22}-s_{12}^2}.
        \end{equation}
        Since $\epsilon>0$,  $S\succ 0$. The following identity holds for a $2\times2$ \glsentryshort{spd} matrix:
        \begin{equation}\label{eq:inv_sqrt_2x2}
        S^{-1/2}
        =
        \sqrt{t+2s}\,(S+s\mathbb{I})^{-1}
        =
        \frac{\sqrt{t+2s}}{2\det(S)+s\,t}
        \begin{pmatrix}
        s_{22}+s & -s_{12}\\
        -s_{12} & s_{11}+s
        \end{pmatrix}.
        \end{equation}
        Therefore, for two modalities the diffusivity matrix is available in closed form, and the Jacobian-space gradient can be evaluated via
        \begin{equation}\label{eq:2x2_gram}
        \nabla_{\mathcal{G}} \phi(\mathcal{G})
        =
        \mathcal{G}\,\Omegaeps{\mathcal{G}},
        \qquad
        \Omegaeps{\mathcal{G}}=S^{-1/2},
        \quad S=G+\epsilon^2 \mathbb{I},
        \end{equation}
        using only dot products and scalar square-roots per voxel.
        
\paragraph{Image-space gradient.}
Define the voxel-wise gradient
\begin{equation}
Q_n
:=
\nabla_{\mathcal{G}}
\phi\!\left(\mathcal{G}_n(\bm{z})\right)
\in \mathbb{R}^{d \times M},
\qquad n=1,\ldots,N,
\end{equation}
and let \(Q = \{Q_n\}_{n=1}^N\).

Since the spatial-gradient field depends linearly on the image through
\begin{equation}
\mathcal{G}(\bm{z}) = D\bm{z},
\end{equation}
the chain rule maps the gradient \(Q\) back to image space using the
adjoint \(D^\ast\):
\begin{equation}
\nabla_{\bm{z}}
\tilde{\mathcal{R}}_{\mathrm{TNV}}(\bm{z})
=
D^\ast Q.
\end{equation}
Defining the discrete divergence by
\begin{equation}
\operatorname{div} := -D^\ast,
\end{equation}
we obtain
\begin{equation}
\nabla_{\bm{z}}
\tilde{\mathcal{R}}_{\mathrm{TNV}}(\bm{z})
=
-\,\operatorname{div} Q,
\end{equation}
where the divergence is applied independently to each column of \(Q\).

\section{Weighted Directional Total Nuclear Variation ($\wdtnv$)}
\label{sec:wdtnv}

This section extends TNV \cite{rigieJointReconstructionMultichannel2015, holtTotalNuclearVariation2014} in two ways:
\begin{itemize}
    \item anatomical information is incorporated via a directional operator
    \item  voxel- and modality-specific diagonal weighting can account for differences in scale and allow each modality to be down-weighted in regions where its measurements are less reliable.

\end{itemize}

\begin{comment}
It is useful to locate this construction within the established taxonomy of multiple-channel regularisation. Channel-separable penalties $\mathcal{R}(\bm{z})=\sum_m \mathcal{R}_0(z_m)$ ignore inter-channel edge correspondence entirely and provide the baseline. Convex multiple-channel penalties couple channels through a mixed norm of the stacked differences, so that an edge present in one channel relaxes the penalty in the others; these are effectively the joint-\gls{tv} family, and correspond to penalising the Frobenius norm of the voxel-wise Jacobian. Rank-based penalties, of which TNV is the canonical example, replace the Frobenius norm by the nuclear norm and so penalise the \emph{number} of distinct edge directions rather than their total magnitude. Line-site and common-sparsity formulations pursue the same goal through auxiliary boundary or support variables, usually at the cost of convexity. The prior developed here sits in the rank-based family and retains its convexity, while drawing anatomical guidance from the separate literature on regularisation with side information, where a co-registered image supplies boundary locations without itself being reconstructed.~\cite[Sections~2.8 and~2.10]{fessler_image_2017}
\end{comment}

\subsection{Directional operators from guidance}
\label{subsec:directional_from_guidance}

We use the directional total-variation operator (Equation~\ref{eq:directional_op}) to incorporate a higher-resolution anatomical image. The image steers the penalty rather than entering the joint Jacobian as an additional column. This keeps the anatomical channel separate from the emission channels, whose units, resolution, and artefacts differ.

If a guidance image were appended as an additional modality, then a single global scaling would have data-dependent and spatially non-uniform consequences. Where emission gradients are weak, even modest guidance gradients could dominate the local Jacobian and hence the singular values; where emission gradients are strong, the same scaling may become negligible. In contrast, a directional operator, $\Pi$, modifies which gradient components are penalised, without forcing the presence of edges where the emission data do not support them (as demonstrated in Chapter~\ref{chapter:deconvolution_pvc}).

Applying a directional operator to each modality gradient defines the directional multi-modality Jacobian
\begin{equation}
\mathcal{G}^{\mathrm{dir}}_n(\bm{z})= \bigl[\Pi_{1,n}(D z_1)_n \;|\; \cdots \;|\; \Pi_{M,n}(D z_M)_n\bigr]\in\mathbb{R}^{d\times M}.
\end{equation}
Here $D$ denotes the chosen discrete directional-difference operator and each $\Pi_{m,n}\in\mathbb{R}^{d\times d}$ is constructed from the guidance field at voxel $n$
\begin{equation}
\xi_n
=
\frac{(D u)_n}{\sqrt{\|(D u)_n\|_2^2+\eta^2}},
\qquad
\Pi_{m,n}=\mathbb{I}_d-\gamma_m\,\xi_n\xi_n^\top,
\qquad \gamma_m\in[0,1].
\label{eq:wdtnv_directional_operator}
\end{equation}
In homogeneous regions of the guidance image, $\xi_n$ is small and $\Pi_{m,n}$ is close to the identity, so the regulariser reduces to an approximately isotropic TNV-type coupling.

Using \gls{ct}  to construct the directional operator  decouples anatomical steering from emission-to-emission weighting. Modality weights determine how strongly emission gradients are coupled to one another through the nuclear norm, while $\Pi_{m,n}$ determines how strongly the penalty prefers guidance-consistent edge directions. Retaining a two-modality Jacobian also preserves the efficient $2\times2$ Gram-matrix structure derived in Section~\ref{sec:tv_to_tnv}.

\subsection{Modality-wise weighting and the full \texorpdfstring{$\wdtnv$}{w-dTNV} prior}
\label{subsec:modality_weighting}

\gls{pet} and \gls{spect} images of the same administered \gls{y90} activity differ substantially in dynamic range: the effective count level, sensitivity, and noise characteristics vary by orders of magnitude between the two modalities. If the directional emission gradients $\Pi_{m,n}(D z_m)_n$ are used directly without normalisation, the nuclear-norm coupling is dominated by the modality with the larger gradient magnitudes.

To compensate, we introduce voxel-wise, modality-wise diagonal weight matrices
\begin{equation}
    B_n := \mathrm{diag}(\rho_{1,n},\,\rho_{2,n}) \in \mathbb{R}^{2\times 2},
    \qquad \rho_{m,n} > 0.
    \label{eq:Bn_def}
\end{equation}
The weighted directional Jacobian is then
\begin{equation}
    \tilde{\mathcal{G}}_n(\bm{z}) := \mathcal{G}_n^{\mathrm{dir}}(\bm{z})\,B_n
    = \bigl[\Pi_{1,n}(D z_1)_n \;\big|\; \Pi_{2,n}(D z_2)_n\bigr]\,
    \mathrm{diag}(\rho_{1,n},\rho_{2,n})
    \in \mathbb{R}^{d\times 2}.
    \label{eq:Yn_weighted}
\end{equation}

Throughout, $B_n$ is treated as a fixed operator that does not depend on $\bm{z}$. In the application of Chapter~\ref{chapter:paper} the weights are computed once, from early-stopped unregularised reconstructions, and then held fixed.

The weighted directional Total Nuclear Variation ($\wdtnv$) prior is then
\begin{equation}
    \mathcal{R}_{\mathrm{w\text{-}dTNV}}(\bm{z})
    :=
    V_{\mathrm{vox}}\sum_{n=1}^{N} \phi\!\left(\tilde{\mathcal{G}}_n(\bm{z})\right),
    \qquad
    \phi(Y) := \sum_{\ell=1}^{\min(d,M)} g(\sigma_\ell(Y)),
    \label{eq:wdtnv_prior}
\end{equation}
where $V_{\mathrm{vox}}$ is the voxel volume and $g$ is the Charbonnier smoothing from~\eqref{eq:charbonnier_g}. Choosing equal weights $\rho_{m,n}\equiv 1$ recovers unweighted $d$TNV. Choosing $\rho_{m,n}$ proportional to the inverse of a characteristic modality scale brings both channels to a comparable range before the nuclear-norm coupling is evaluated. The \gls{pet}/\gls{spect} normalisation used in the application study is specified in Chapter~\ref{chapter:paper}.

\section{Conclusion}

This chapter introduced a smoothed, weighted, and directional Total Nuclear Variation ($\wdtnv$) prior for synergistic reconstruction. Charbonnier smoothing of the singular values results in a differentiable and convex spectral functional whose Jacobian-space gradient can be written in Gram-matrix form. Convexity follows from absolute symmetry and is inherited by the full prior because $D$, $\Pi_{m,n}$ and $B_n$ are all fixed linear maps. For two coupled emission modalities the key spectral operation reduces to the analytical calculation of the inverse square-root of   $2\times2$ \gls{spd} matrices, allowing the prior and its gradient to be evaluated robustly without a general per-voxel \glsentryshort{svd}.

The chapter also separated two roles that are easily confounded in multimodality reconstruction. Anatomical information is used directionally through $\Pi_{m,n}$, steering the penalty without adding \gls{ct} as another coupled intensity channel. Modality scaling is handled separately through the diagonal weight matrix $B_n$, allowing \gls{pet} and \gls{spect} gradients to enter the nuclear-norm coupling on comparable scales. Chapter~\ref{chapter:optimisation} uses these definitions to construct practical preconditioners and subset-selection strategies for large-scale $\wdtnv$ optimisation.
